\section{High-Dimensional Effects in k-NN}

\begin{frame}{Curse of Dimensionality for k-NN}
\begin{itemize}
  \item As dimension grows, points become similarly distant.
  \pause
  \item Nearest and farthest neighbors become less distinguishable.
  \pause
  \item Local neighborhoods become less local.
  \pause
  \item Brute-force search also becomes expensive with high $d$ and large $n$.
\end{itemize}

\pause
\begin{rlintuition}
k-NN is strongest when representation space has meaningful geometry.
\end{rlintuition}
\end{frame}

\begin{frame}{Graphic: Distance Concentration in High Dimensions}
\centering
\includegraphics[width=0.84\textwidth]{figures/knn/knn_distance_concentration.pdf}

\vspace{0.35em}
\footnotesize Blue line increasing to 1 means nearest and farthest distances become closer.
\end{frame}

\begin{frame}{Making k-NN Work in High Dimensions}
\begin{itemize}
  \item Standardize or normalize features before distance computation.
  \pause
  \item Reduce dimension (PCA, feature selection, domain features).
  \pause
  \item Use better representations (contrastive/self-supervised embeddings).
  \pause
  \item Choose metric for representation type (cosine for embeddings is common).
  \pause
  \item Evaluate robustness under feature noise and distribution shift.
\end{itemize}
\end{frame}
